\documentclass{article}
\usepackage[margin=1in]{geometry}
\usepackage[skip=1em, indent=12pt]{parskip}
\usepackage{amsmath}
\usepackage{amssymb}
\usepackage{graphicx}
\usepackage{microtype}
\usepackage{textcomp, gensymb}
\usepackage{enumitem}
\begin{document}
\setcounter{section}{4}
\section{Discrete-Time System Analysis using the z-Transform} 
\setcounter{subsection}{0}
\subsection{The z-Transform}
Put simply, the z-Transform is the discrete-time equivalent of the Laplace Transform, defined mathematically as
\[
	X[z] = \sum_{n=-\infty}^{\infty}x[n]z^{-n},
\]
where \(z = re^{j\theta}\). However, we'll never need to use this formula (or its inverse) specifically, and will refer to a table of transform pairs in most cases.

Similar to the Laplace Transform, the z-Transform is \textbf{linear}, meaning they obey the following property:
\[
	a_1x_1[n] + a_2x_2[n] \Leftrightarrow a_1X_1[z] + a_2X_2[z].
\]
The unilateral z-Transform, defined as
\[
	X[z] = \sum_{n=0}^{\infty}x[n]z^{-n},
\]
is going to be used most often here, since the signals we will be dealing with are largely causal, and the bilateral transform is not needed. If we \textit{do} use the bilateral transform, the region of convergence of the z-Transform will become relevant.

\subsubsection{Example:}

Find the z-transform and the corresponding ROC for the signal \(\gamma^nu[n]\). 

First, recall that the given function represents a signal that grows or decays exponentially (depending on the value of \(\gamma\)), beginning at \(n=0\), where the signal's value \(x[0] = 1\), and either decays or grows for \(n>0\).

Formulating the transform mathematically, 
\[
	X[z] = \sum_{n=0}^{\infty}\gamma^nu[n]z^{-n}.
\]

Since the unit step function will always be equal to \(u[n] = 1\) for all values of \(n \geq 0\), we may simplify the transform to 
\[
	X[z] = \sum_{n=0}^{\infty}\left(\frac{\gamma}{z}\right)^n,
\]
which reduces to a geometric series of the form 
\[
  1 + x + x^2 + x^3 + ... = \frac{1}{1-x},
\]
if \(x < 1\).

Therefore, we may simplify the infinite sum produced by the z-transform formula to
\[
	X[z] = \frac{1}{1-\frac{\gamma}{2}} = \frac{z}{z-\gamma}
\]
as long as \(|z| > |\gamma|\). This condition defines the \textit{region of convergence} for the bilateral transform.

When using the unilateral z-transform, we will follow a very similar process to our use of the Laplace transform thus far, matching a given function to a transform pair or several transform pairs in the given table, potentially making changes to the original function to match it more closely to the table.

For the inverse z-transform, we'll once again use partial fraction expansion.

\setcounter{subsubsection}{2}
\subsubsection{Example:}
Find the inverse z-transform of 
\[
  \frac{8z-19}{(z-2)(z-3)}.
\]

Here, we'll expand the given function to
\[
	X[z] = \frac{3}{z-2} + \frac{5}{z-3}
\]
and use the transform pair
\[
	\gamma^{n-1}u[n-1] \Leftrightarrow \frac{1}{z-\gamma}
\]
to find the solution as
\[
	x[n] = [3(2)^{n-1} + 5(3)^{n-1}]u[n-1].
\]

To get a function using \(u[n]\) instead, we may use \textbf{modified} partial fraction expansion, using \(X[z]/z\) instead:
\[
	\frac{X[z]}{z} = \frac{-19/6}{z} + \frac{3/2}{z-2} + \frac{5/3}{z-3}.
\]
Using the transform pair
\[
	\gamma^nu[n] \Leftrightarrow \frac{z}{z-\gamma},
\]
we may find the solution as 
\[
	x[n] = -\frac{19}{6}\delta[n] + [\frac{3}{2}(2)^n + \frac{5}{3}(3)^n]u[n].
\]

Another important correlation between the Laplace- and z-Transform is the \textbf{system transfer function}, defined as
\[
	H[z] = \sum_{n=-\infty}^{\infty}h[n]z^{-n}
\]
and
\[
	h[n] \Leftrightarrow H[z].
\]




\end{document}
